\section{结论}
\label{sec:conclusion}

我们提出了\fasa{}（\fasafull{}），一个通过帧间特征时间对比自监督发现任务相关空间区域的即插即用空间注意力模块。\fasa{}的核心洞察简单而强大：连续帧之间变化最显著的区域对应于当前活跃的交互区域。基于这一洞察，\fasa{}将空间特征分解为动态流（交互区域）和静态流（稳定上下文），编码变化质心和运动方向，并提取空间关键点坐标——提供比全局平均池化丰富$2.07\times$的空间信息，每个相机仅增加约67k参数。

\fasa{}完全自监督，无需空间标注、外部教师模型或专门的辅助损失。它仅需连续观测帧并在标准卷积特征图上操作，使其广泛适用于本文研究的特定架构和任务之外。结合特征差分编码的"语义速度"，\fasa{}为流匹配策略提供了丰富的空间感知和隐式的时间变化感知。

我们将\fasa{}集成到VITA流匹配架构中，在具有挑战性的6自由度S型焊缝跟踪任务上进行评估。实验表明\fasa{}显著提升了跟踪精度，消融实验证实了各空间特征组件（动态池化、静态池化、变化质心编码、空间softmax坐标）的互补贡献。

我们认为自监督空间注意力是机器人模仿学习中广泛重要的能力。\fasa{}所采取的方法——从观测流本身提取空间注意力，而非依赖人工标注或外部教师模型——代表了构建更强大、更便携、更具泛化能力的机器人策略的有前景的方向。
